% =====================================================================
%  Thesis appendix: Sequential Linear Programming -- Implementation
%  --------------------------------------------------------------------
%  Expands the compact SLP description given in the main thesis section
%  \ref{sec:SLP}. The appendix documents the implemented merit-function
%  SLP solver, inspired by the MATLAB framework of \cite{sorensen2026fminslp},
%  and its coupling to the MAPDL beam-model pipeline.
%
%  Required preamble packages:
%      amsmath, amssymb, booktabs, enumitem, hyperref, algorithm,
%      algpseudocode
%
%  If \citet is used instead of \cite, add natbib or use the citation
%  commands from the main thesis template.
% =====================================================================

\documentclass[11pt,a4paper]{article}

\usepackage[margin=25mm]{geometry}
\usepackage{amsmath}
\usepackage{amssymb}
\usepackage{booktabs}
\usepackage{enumitem}
\usepackage{hyperref}
\usepackage{algorithm}
\usepackage{algpseudocode}

\title{Sequential Linear Programming -- Implementation Details\\
       \large Appendix supporting \autoref{sec:SLP}}
\author{DMS4 Optimization Appendix}
\date{\today}

\begin{document}

\maketitle

% =====================================================================
\section{Sequential Linear Programming -- Implementation Details}
\label{sec:SLP_appendix}

This appendix expands the Sequential Linear Programming (SLP)
description given in \autoref{sec:SLP}. The purpose is to
document the numerical procedure used to optimize the LT1292 MAPDL
beam model, including the linearized subproblem, slack relaxation,
merit-function acceptance, filter acceptance, move-limit adaptation,
penalty update and stopping criteria.

The implemented optimizer is a Python implementation inspired by the
merit-function variant of the MATLAB framework described by
\cite{sorensen2026fminslp}. The solver is coupled to the MAPDL
analysis pipeline described in the code-implementation chapter. In
each outer iteration, the current structural design is evaluated in
MAPDL, local finite-difference sensitivities are computed, and a
linear programming (LP) subproblem is solved inside a move-limit
trust region. The resulting trial design is then checked by a new
MAPDL evaluation before it is accepted or rejected.

The appendix is structured as follows. \autoref{sec:slpA:problem}
defines the optimization problem and the sign convention used by the
implementation. \autoref{sec:slpA:linearise} describes the local
linearization and the finite-difference gradient calculation.
\autoref{sec:slpA:lp} presents the slack-relaxed LP subproblem.
\autoref{sec:slpA:merit} introduces the merit function used both
in the LP objective and in the nonlinear acceptance test.
\autoref{sec:slpA:accept} explains the filter, sufficient-decrease
test and feasibility-restoration logic. \autoref{sec:slpA:movelim}
describes the adaptive move limits, which act as a component-wise
trust region. \autoref{sec:slpA:penalty} and
\autoref{sec:slpA:termination} describe the adaptive slack penalty and
the termination criteria. The full iteration is summarized in
\autoref{alg:slp}, and the numerical settings used in the present
study are listed in \autoref{sec:slpA:settings}.

% =====================================================================
\subsection{Problem statement and sign convention}
\label{sec:slpA:problem}

The structural optimization problem is formulated as a constrained
mass-minimization problem. The design vector is denoted
$x \in \mathbb{R}^n$, where each component represents a physical
design variable in the MAPDL beam model, such as a tube radius,
diameter or wall thickness. The problem is written as
\begin{equation}
\label{eq:appA_problem}
\begin{aligned}
    \min_{x \in \mathbb{R}^n} \quad & f(x), \\
    \text{s.t.} \quad
    & c_i(x) \geq 0,
    \qquad i = 1,\dots,m, \\
    & x^L \leq x \leq x^U .
\end{aligned}
\end{equation}
Here, $f(x)$ is the structural mass returned by the MAPDL analysis.
The vector $c(x)$ collects the design constraints, including
utilization constraints, cross-section class constraints,
eigenvalue constraints and segment-mass constraints. The lower and
upper bounds $x^L$ and $x^U$ define the physically admissible design
space.

In the thesis formulation, the constraints are written in the
non-negative convention $c_i(x) \geq 0$. In this convention, a
positive value means that the constraint is satisfied, while a
negative value means that the constraint is violated. The SLP
implementation uses the opposite convention,
\begin{equation}
\label{eq:appA_g_convention}
    g_i(x) \leq 0,
    \qquad i = 1,\dots,m .
\end{equation}
The two formulations are connected by a sign change in the constraint
callback,
\begin{equation}
\label{eq:appA_sign_flip}
    g(x) = -c(x).
\end{equation}
All algorithmic quantities in the following sections are therefore
defined using the $g(x) \leq 0$ convention. Thus, a constraint is
satisfied when $g_i(x) \leq 0$ and violated when $g_i(x)>0$.

% =====================================================================
\subsection{Local linearization and finite-difference gradients}
\label{sec:slpA:linearise}

Sequential Linear Programming solves a nonlinear constrained problem
by repeatedly replacing it with a local linear approximation. At
iteration $k$, the current design is $x_k$, and the trial design is
written as
\begin{equation}
\label{eq:appA_trial_design}
    x_{\text{trial}} = x_k + \Delta x ,
\end{equation}
where $\Delta x$ is the design change proposed by the LP subproblem.

The objective function is approximated by the first-order Taylor
expansion
\begin{equation}
\label{eq:appA_lin_f}
    f(x_k + \Delta x)
    \approx
    f(x_k) + \nabla f(x_k)^{\!\top}\Delta x .
\end{equation}
Similarly, the constraint vector is approximated as
\begin{equation}
\label{eq:appA_lin_g}
    g(x_k + \Delta x)
    \approx
    g(x_k) + J(x_k)\Delta x .
\end{equation}
The Jacobian $J(x_k) \in \mathbb{R}^{m \times n}$ contains the local
constraint sensitivities,
\begin{equation}
\label{eq:appA_jacobian}
    J_{ij}(x_k)
    =
    \frac{\partial g_i(x_k)}{\partial x_j},
    \qquad
    i = 1,\dots,m,\quad j = 1,\dots,n .
\end{equation}

The MAPDL model is treated as a black-box analysis model. Analytical
gradients are therefore not available. Instead, the objective
gradient and the constraint Jacobian are estimated by forward finite
differences. For design variable $i$, the perturbation size is chosen
as
\begin{equation}
\label{eq:appA_fd_step}
    h_i =
    h_{\mathrm{rel}}
    \max\left(1, |x_{k,i}|\right),
\end{equation}
where $h_{\mathrm{rel}}$ is a fixed relative finite-difference step.
Using the unit vector $e_i$, the objective sensitivity is estimated
as
\begin{equation}
\label{eq:appA_fd_f}
    \frac{\partial f(x_k)}{\partial x_i}
    \approx
    \frac{
        f(x_k + h_i e_i) - f(x_k)
    }{h_i},
\end{equation}
and the corresponding constraint sensitivities are estimated as
\begin{equation}
\label{eq:appA_fd_g}
    \frac{\partial g_j(x_k)}{\partial x_i}
    \approx
    \frac{
        g_j(x_k + h_i e_i) - g_j(x_k)
    }{h_i},
    \qquad j = 1,\dots,m .
\end{equation}

For $n$ design variables, one gradient evaluation requires one
baseline MAPDL evaluation and $n$ perturbed evaluations. The finite
difference calculation is therefore the dominant computational cost
of each outer SLP iteration.

% =====================================================================
\subsection{Slack-relaxed linear subproblem}
\label{sec:slpA:lp}

If the local linear models in \autoref{eq:appA_lin_f} and
\autoref{eq:appA_lin_g} are used directly, the SLP subproblem becomes
\begin{equation}
\label{eq:appA_direct_lp}
\begin{aligned}
    \min_{\Delta x}\quad
        & \nabla f(x_k)^{\!\top}\Delta x ,
    \\
    \text{s.t.}\quad
        & g(x_k) + J(x_k)\Delta x \leq 0 ,
    \\
        & x^{L,\mathrm{cur}}
          \leq
          x_k + \Delta x
          \leq
          x^{U,\mathrm{cur}} .
\end{aligned}
\end{equation}
The constant term $f(x_k)$ is omitted from the objective because it
does not affect the minimizer of the LP.

The direct LP in \autoref{eq:appA_direct_lp} may become infeasible if
the current design violates one or more constraints and the
move-limit box is too small to recover feasibility in one step. This
situation is expected in the present application because one
LT1292 design may violate some stability-related constraints. To make
the LP solvable from infeasible points, each constraint is
relaxed by a non-negative slack variable $y_i$,
\begin{equation}
\label{eq:appA_slack}
    g_i(x) - y_i \leq 0,
    \qquad
    y_i \geq 0,
    \qquad
    i = 1,\dots,m .
\end{equation}
The slack variable represents the amount of artificial constraint
relaxation. If $y_i = 0$, the original constraint is enforced. If
$y_i > 0$, the corresponding constraint is temporarily relaxed.
Because permanent slack usage would correspond to an infeasible
engineering design, the slacks are penalized through the merit
function introduced in \autoref{sec:slpA:merit}.

The initial slack vector is chosen as
\begin{equation}
\label{eq:appA_slack_init}
    y_i^0 = \max\left(0, g_i(x_0)\right),
    \qquad i = 1,\dots,m ,
\end{equation}
so that the initial slack values cover the initial constraint
violations. To keep the LP bounded, the slacks are also assigned a
large finite upper bound,
\begin{equation}
\label{eq:appA_slack_upper}
    y^U =
    10^6 \max\left(1,\|g(x_0)\|_\infty\right).
\end{equation}
This upper bound is chosen large enough that it does not normally
restrict the solution, but it prevents the LP from becoming
unbounded in pathological cases.

The augmented variable and its increment are defined as
\begin{equation}
\label{eq:appA_augmented_variable}
    z =
    \begin{bmatrix}
        x \\ y
    \end{bmatrix},
    \qquad
    \Delta z =
    \begin{bmatrix}
        \Delta x \\ \Delta y
    \end{bmatrix}.
\end{equation}
The slack-relaxed LP solved at each outer iteration is then
\begin{equation}
\label{eq:appA_slack_relaxed_lp}
\begin{aligned}
    \min_{\Delta z}\quad
        & \nabla_z \Phi(x_k,y_k)^{\!\top}\Delta z ,
    \\
    \text{s.t.}\quad
        & g(x_k) + J(x_k)\Delta x - (y_k + \Delta y) \leq 0 ,
    \\
        & x^{L,\mathrm{cur}}
          \leq
          x_k + \Delta x
          \leq
          x^{U,\mathrm{cur}},
    \\
        & 0
          \leq
          y_k + \Delta y
          \leq
          y^U .
\end{aligned}
\end{equation}
The current bounds $x^{L,\mathrm{cur}}$ and $x^{U,\mathrm{cur}}$ are
the move-limit bounds described in \autoref{sec:slpA:movelim}.

The linearized relaxed constraint in \autoref{eq:appA_slack_relaxed_lp}
can be written in the standard LP form
\begin{equation}
\label{eq:appA_lp_standard_form}
    A_{\mathrm{ub}}\Delta z \leq b_{\mathrm{ub}},
\end{equation}
where
\begin{equation}
\label{eq:appA_lp_standard_matrices}
    A_{\mathrm{ub}}
    =
    \begin{bmatrix}
        J(x_k) & -I_m
    \end{bmatrix},
    \qquad
    b_{\mathrm{ub}}
    =
    -\left(g(x_k)-y_k\right).
\end{equation}
Because each inequality has its own non-negative slack variable with
a large finite upper bound, the LP remains practically feasible for
the constraint violations encountered in the present study. The LP is
solved using a standard linear programming routine.

% =====================================================================
\subsection{Merit function with linear-quadratic slack penalty}
\label{sec:slpA:merit}

The slack variables make the LP robust against infeasible iterates,
but they must not remain active in the final design. The algorithm
therefore uses a scalar merit function that combines the objective
value and the slack violation into one quantity. The merit function
is defined as
\begin{equation}
\label{eq:appA_merit}
    \Phi(x,y)
    =
    f(x)
    +
    a_F
    \sum_{i=1}^{m}
    \left(
        R y_i
        +
        \frac{1}{2}y_i^2
    \right),
\end{equation}
where $R > 0$ is the slack penalty and
\begin{equation}
\label{eq:appA_af}
    a_F = \max\left(|f(x_0)|,1\right)
\end{equation}
is an objective scaling factor fixed at the first MAPDL evaluation.

The merit function in \autoref{eq:appA_merit} has two roles. First,
it discourages the optimizer from using slacks permanently. Second,
it provides a scalar quantity that can be used to compare the
predicted LP improvement with the actual nonlinear improvement after
a trial design has been evaluated in MAPDL.

The slack penalty contains both a linear term and a quadratic term.
The linear term $R y_i$ makes even small slack values expensive,
while the quadratic term $\frac{1}{2}y_i^2$ increasingly penalizes
larger violations. The scaling factor $a_F$ makes the slack penalty
less sensitive to the units and magnitude of the mass objective.

The gradient of the merit function with respect to the augmented
variable $z=[x,y]$ is
\begin{equation}
\label{eq:appA_merit_grad}
    \nabla_z \Phi(x,y)
    =
    \begin{bmatrix}
        \nabla f(x)
        \\
        a_F(R\mathbf{1}+y)
    \end{bmatrix},
\end{equation}
where $\mathbf{1}$ is a vector of ones of length $m$. This gradient is
used as the cost vector in the LP objective in
\autoref{eq:appA_slack_relaxed_lp}. The LP therefore seeks a step that
reduces the linearized merit function, not only the structural mass.
This couples the LP subproblem to the nonlinear acceptance test
described in \autoref{sec:slpA:accept}.

% =====================================================================
\subsection{Step acceptance}
\label{sec:slpA:accept}

The LP solution provides a candidate step
$\Delta z=[\Delta x,\Delta y]$, but the LP itself does not guarantee
that the step is good for the original nonlinear MAPDL model. This is
because the LP is based only on local linear approximations. The
trial design must therefore be evaluated in MAPDL and accepted only
if it gives sufficient progress.

The implementation uses three acceptance mechanisms: a filter test,
a merit-function sufficient-decrease test and a feasibility
restoration test.

% ---------------------------------------------------------------------
\paragraph{Filter acceptance.}

The filter compares trial points using two quantities: a measure of
constraint violation and a scaled merit value. The relaxed maximum
constraint violation is defined as
\begin{equation}
\label{eq:appA_filter_h}
    h(x,y)
    =
    \max_i
    \max\left(0, g_i(x)-y_i\right),
\end{equation}
and the scaled merit value is
\begin{equation}
\label{eq:appA_filter_f}
    f_{\mathrm{filter}}(x,y)
    =
    \frac{\Phi(x,y)}{\Phi_0},
    \qquad
    \Phi_0=\Phi(x_0,y_0).
\end{equation}
A trial point is represented by the pair
\begin{equation}
\label{eq:appA_filter_pair}
    \left(h^*, f^*\right)
    =
    \left(
        h(x_{\mathrm{trial}},y_{\mathrm{trial}}),
        f_{\mathrm{filter}}(x_{\mathrm{trial}},y_{\mathrm{trial}})
    \right).
\end{equation}

Following the filter concept used in \cite{sorensen2015thicknessfilters},
a trial point is accepted by the filter $\mathcal{F}_k$ if it is not
dominated by any stored filter entry. In the implemented test, this
is written as
\begin{equation}
\label{eq:appA_filter_accept}
    h^* \leq \beta h_j
    \quad \text{or} \quad
    f^* + \gamma h^* \leq f_j,
    \qquad
    \forall (h_j,f_j)\in\mathcal{F}_k ,
\end{equation}
where $\beta,\gamma\in(0,1)$ are small filter parameters.

This condition means that a trial point may be accepted if it gives
sufficiently lower constraint violation or sufficiently lower scaled
merit value compared with all stored filter entries. In this way, the
filter prevents the algorithm from accepting points that are worse in
both feasibility and merit value. After an accepted infeasible step,
the new filter pair is inserted into the filter and entries dominated
by the new pair are removed. Feasible points with $h^*=0$ are not
stored.

% ---------------------------------------------------------------------
\paragraph{Sufficient decrease in the merit function.}

Filter acceptance alone is not sufficient, because the filter does
not directly compare the actual nonlinear reduction with the
reduction predicted by the LP model. The implementation therefore
also checks a merit-function decrease condition.

The actual merit decrease is
\begin{equation}
\label{eq:appA_merit_actual}
    \Delta \Phi_{\mathrm{act}}
    =
    \Phi(x_k,y_k)
    -
    \Phi(x_k+\Delta x, y_k+\Delta y),
\end{equation}
while the predicted merit decrease from the linear LP model is
\begin{equation}
\label{eq:appA_merit_pred}
    \Delta \Phi_{\mathrm{pred}}
    =
    -
    \nabla_z\Phi(x_k,y_k)^{\!\top}\Delta z .
\end{equation}
If the LP predicts a merit decrease,
$\Delta\Phi_{\mathrm{pred}}>0$, the accepted trial point must satisfy
the sufficient-decrease condition
\begin{equation}
\label{eq:appA_armijo}
    \Delta \Phi_{\mathrm{act}}
    \geq
    \sigma \Delta \Phi_{\mathrm{pred}},
    \qquad
    \sigma\in(0,1).
\end{equation}
This condition requires the actual nonlinear merit decrease to be at
least a small fraction of the decrease predicted by the linear model.

If $\Delta\Phi_{\mathrm{pred}}\leq 0$, the LP step is not a predicted
descent step for the merit function. In this case, the sufficient
decrease test is not informative as a merit-decrease criterion, and
the step must instead be justified by the filter test or the
feasibility restoration test.

% ---------------------------------------------------------------------
\paragraph{Feasibility restoration.}

During early iterations, the current design may be infeasible with
respect to the original physical constraints. A trial step that fails
the filter or merit-decrease test may still be useful if it moves the
design significantly closer to feasibility. The total original
constraint violation is measured as
\begin{equation}
\label{eq:appA_total_violation}
    V(x)
    =
    \sum_{i=1}^{m}
    \max\left(0,g_i(x)\right).
\end{equation}
If the current design is infeasible,
$V(x_k)>\varepsilon_{\mathrm{feas}}$, a trial point may be accepted
by feasibility restoration if
\begin{equation}
\label{eq:appA_feas}
    V(x_{\mathrm{trial}})
    \leq
    (1-\gamma_v)V(x_k),
    \qquad
    \gamma_v\in(0,1).
\end{equation}
This criterion allows the optimizer to prioritize feasibility
improvement when the initial design violates several constraints.

% ---------------------------------------------------------------------
\paragraph{Trust-region backtracking.}

If neither the merit/filter acceptance route nor the feasibility
restoration route accepts the trial point, the implementation does
not perform a classical fractional line search along the same
direction. Instead, the move-limit box is contracted and the LP is
solved again inside the smaller trust region. This is more accurately
described as trust-region backtracking than as line-search
backtracking.

The contraction is performed using the factor
$\eta_{\mathrm{bt}}\in(0,1)$. If the current move-limit vector is
denoted by $\Delta x^{\max}$, a rejected trial updates it according to
\begin{equation}
\label{eq:appA_bt_shrink}
    \Delta x^{\max}
    \leftarrow
    \eta_{\mathrm{bt}}\Delta x^{\max}.
\end{equation}
The LP is then re-solved with the reduced move-limit box. At most
$N_{\mathrm{bt}}$ such inner attempts are allowed per outer SLP
iteration.

% =====================================================================
\subsection{Adaptive move limits as a trust region}
\label{sec:slpA:movelim}

The local linear approximation is only reliable near the current
design. The SLP method therefore restricts each LP step to a
component-wise move-limit box,
\begin{equation}
\label{eq:appA_move_box}
    x^{L,\mathrm{cur}}
    \leq
    x_k + \Delta x
    \leq
    x^{U,\mathrm{cur}} .
\end{equation}
This box acts as a trust region. Large steps are avoided because they
would make the linearized objective and constraints unreliable.

The move limits are updated separately for each design variable. This
is important for the LT1292 because the design variables have
different physical scales. For example, chord diameters may be of
order $60~\mathrm{mm}$, while brace wall thicknesses may be of order
$2~\mathrm{mm}$. A single scalar move limit in physical units would
therefore be inappropriate.

Let \autoref{eq:appA_step_history}
\begin{equation}
\label{eq:appA_step_history}
    d_{i,k}
    =
    x_{i,k} - x_{i,k-1}
\end{equation}
be the previous accepted change in variable $i$. The sign consistency
between two consecutive accepted changes is measured by
\begin{equation}
\label{eq:appA_sign_product}
    q_{i,k}
    =
    d_{i,k}d_{i,k-1}.
\end{equation}
If $q_{i,k}<0$, the variable has changed direction between
iterations, indicating oscillatory behavior. The move limit for that
variable is then reduced by
\begin{equation}
\label{eq:appA_move_shrink}
    \Delta x_i^{\max}
    \leftarrow
    a_{\mathrm{shr}}\Delta x_i^{\max},
    \qquad
    0<a_{\mathrm{shr}}<1 .
\end{equation}
If $q_{i,k}>0$, the variable has moved in the same direction for two
consecutive iterations. This indicates coherent progress, and the
move limit is expanded by
\begin{equation}
\label{eq:appA_move_expand}
    \Delta x_i^{\max}
    \leftarrow
    a_{\mathrm{exp}}\Delta x_i^{\max},
    \qquad
    a_{\mathrm{exp}}>1 .
\end{equation}
If one of the two previous changes is numerically negligible, the
history is not considered reliable and the move limit is reset to the
base value.

After this update, the move limits are clipped between a lower and an
upper value,
\begin{equation}
\label{eq:appA_move_bounds}
    \delta_{\min}(x_i^U-x_i^L)
    \leq
    \Delta x_i^{\max}
    \leq
    \delta_{\max}(x_i^U-x_i^L).
\end{equation}
The current LP bounds are then formed as
\begin{equation}
\label{eq:appA_current_bounds}
\begin{aligned}
    x_i^{L,\mathrm{cur}}
    &=
    \max\left(x_i^L, x_{i,k}-\Delta x_i^{\max}\right),
    \\
    x_i^{U,\mathrm{cur}}
    &=
    \min\left(x_i^U, x_{i,k}+\Delta x_i^{\max}\right).
\end{aligned}
\end{equation}

Thus, variables that oscillate are treated more conservatively,
whereas variables that move consistently are allowed to take larger
steps. This adaptive rule is a simple trust-region mechanism tailored
to SLP with bound-constrained design variables.

% =====================================================================
\subsection{Adaptive slack penalty}
\label{sec:slpA:penalty}

The slack penalty $R$ controls how expensive it is for the optimizer
to use artificial constraint relaxation. It appears both in the LP
cost vector through \autoref{eq:appA_merit_grad} and in the nonlinear
merit function in \autoref{eq:appA_merit}. A small value of $R$
allows the optimizer to use slacks easily, while a large value forces
the optimizer to reduce constraint violations more aggressively.

A fixed value of $R$ is not robust for all stages of the optimization.
Early in the run, some slack may be necessary to move from an
infeasible starting design. Later in the run, however, slacks should
be driven to zero. The implementation therefore increases $R$
monotonically when the optimization shows signs of relying too much
on slack variables. The update is
\begin{equation}
\label{eq:appA_penalty}
    R_{k+1}
    =
    \min\left(a_M R_k, R_{\max}\right),
    \qquad
    a_M>1 .
\end{equation}
The penalty is increased if at least one of the following conditions
is met:
\begin{enumerate}[label=(\roman*)]
    \item the LP slacks remain active,
    \begin{equation}
    \label{eq:appA_slack_active}
        \sum_{i=1}^m y_i > \varepsilon_{\mathrm{slack}},
    \end{equation}

    \item the original feasibility violation stalls,
    \begin{equation}
    \label{eq:appA_feas_stall}
        V(x_{k+1}) \geq 0.95 V(x_k),
    \end{equation}

    \item all inner LP trial points are rejected.
\end{enumerate}
The upper bound $R_{\max}$ prevents the slack term from becoming
numerically dominant in pathological cases.

% =====================================================================
\subsection{Termination criteria}
\label{sec:slpA:termination}

The solver reports success only at an accepted iterate that satisfies
the original constraints to the prescribed feasibility tolerance.
This requirement is important because a design with nonzero slack is
not an acceptable engineering design.

Four numerical stopping indicators are evaluated. The first is an
objective-gradient indicator,
\begin{equation}
\label{eq:appA_stop_grad}
    \|\nabla f(x_k)\|_2
    \leq
    \varepsilon_g .
\end{equation}
For a constrained problem, this is not a full Karush--Kuhn--Tucker
optimality check, since the active constraint gradients and Lagrange
multipliers are not explicitly included. It is therefore used as a
practical stationarity indicator rather than as a rigorous constrained
optimality condition.

The second stopping indicator is the merit-change tolerance,
\begin{equation}
\label{eq:appA_stop_merit}
    \left|
        \Phi(x_{k+1},y_{k+1})
        -
        \Phi(x_k,y_k)
    \right|
    \leq
    \varepsilon_f .
\end{equation}
The third stopping indicator is the accepted step length,
\begin{equation}
\label{eq:appA_stop_step}
    \|x_{k+1}-x_k\|_2
    \leq
    \varepsilon_x .
\end{equation}
The fourth stopping indicator is a lower bound on the merit value,
\begin{equation}
\label{eq:appA_stop_obj_limit}
    \Phi(x_{k+1},y_{k+1})
    \leq
    \Phi_{\min},
\end{equation}
which allows the solver to terminate early if the merit drops below a
user-specified lower bound $\Phi_{\min}$. In the present study,
$\Phi_{\min}$ is set sufficiently low that this indicator never
triggers; it is included for completeness because the implementation
checks it on every accepted trial.

A successful termination requires that at least one of the four
stopping indicators is satisfied and that the accepted design is
feasible:
\begin{equation}
\label{eq:appA_stop_feas}
    \max_i \max\left(0,g_i(x_{k+1})\right)
    \leq
    \varepsilon_{\mathrm{feas}},
\end{equation}
with negligible remaining slack,
\begin{equation}
\label{eq:appA_stop_slack}
    \max_i y_{i,k+1}
    \leq
    \max\left(
        \varepsilon_{\mathrm{slack}},
        \varepsilon_{\mathrm{feas}}
    \right).
\end{equation}
The maximum number of outer iterations is limited by $N_{\max}$. In
addition, the total number of MAPDL evaluations is bounded by
\begin{equation}
\label{eq:appA_stop_nfev}
    N_{\mathrm{fev}} \leq 10\,N_{\max}(n+1),
\end{equation}
where $N_{\mathrm{fev}}$ is the running count of baseline,
finite-difference and inner-trial evaluations. This budget acts as a
hard ceiling on the total computational cost and prevents runaway
runs in which $N_{\max}$ outer iterations are never reached because
each iteration consumes many inner trials.

If the move-limit box collapses to the lower bound
$\delta_{\min}$ without satisfying the feasibility and stopping
requirements, the solver reports failure. In this context, failure
does not necessarily mean that no better design exists. It usually
indicates that the local linear model is no longer producing reliable
accepted steps. Possible causes include finite-difference noise,
poorly scaled constraints, a strongly nonlinear active constraint or
an overly restrictive move-limit box.

% =====================================================================
\subsection{Algorithm summary}
\label{sec:slpA:algo}

\autoref{alg:slp} summarizes one outer iteration of the implemented
SLP solver.

\begin{algorithm}[h]
\caption{One outer iteration of the implemented SLP solver.}
\label{alg:slp}
\begin{algorithmic}[1]
\State Evaluate $f(x_k)$ and $g(x_k)=-c(x_k)$ using the MAPDL model.
\State Estimate $\nabla f(x_k)$ and $J(x_k)$ by forward finite differences.
\State Form the merit gradient $\nabla_z\Phi(x_k,y_k)$ using \autoref{eq:appA_merit_grad}.
\State Update the component-wise move-limit box using \autoref{eq:appA_current_bounds}.
\State Set \texttt{accepted} $\gets$ \textbf{false}.
\For{$j=1,\dots,N_{\mathrm{bt}}$}
    \State Solve the slack-relaxed LP in \autoref{eq:appA_slack_relaxed_lp} for $(\Delta x,\Delta y)$.
    \State Set $x_{\mathrm{trial}} \gets x_k+\Delta x$.
    \State Set $y_{\mathrm{trial}} \gets y_k+\Delta y$.
    \State Evaluate $f(x_{\mathrm{trial}})$ and $g(x_{\mathrm{trial}})$ using MAPDL.
    \State Compute the filter pair using \autoref{eq:appA_filter_pair}.
    \State Compute $\Delta\Phi_{\mathrm{act}}$ and $\Delta\Phi_{\mathrm{pred}}$ using \autoref{eq:appA_merit_actual} and \autoref{eq:appA_merit_pred}.
    \If{filter acceptance in \autoref{eq:appA_filter_accept} is satisfied
        \textbf{and} (sufficient decrease in \autoref{eq:appA_armijo}
        \textbf{or} feasibility restoration in \autoref{eq:appA_feas} holds)}
        \State Accept trial:
        $x_{k+1}\gets x_{\mathrm{trial}}$,
        $y_{k+1}\gets y_{\mathrm{trial}}$.
        \State Update the filter if required.
        \State Set \texttt{accepted} $\gets$ \textbf{true}.
        \State \textbf{break}.
    \Else
        \State Contract the move-limit box using \autoref{eq:appA_bt_shrink}.
    \EndIf
\EndFor
\If{\texttt{accepted} is \textbf{false}}
    \State Increase the slack penalty if required.
    \State Continue to the next outer iteration or report failure if the move limits have collapsed.
\EndIf
\If{slacks are active or feasibility has stalled}
    \State Update $R \gets \min(a_MR,R_{\max})$ using \autoref{eq:appA_penalty}.
\EndIf
\If{the accepted iterate satisfies the feasibility and stopping criteria in \autoref{eq:appA_stop_feas}--\autoref{eq:appA_stop_slack}}
    \State \Return \textbf{success}.
\EndIf
\If{the maximum number of outer iterations has been reached}
    \State \Return \textbf{failure}.
\EndIf
\end{algorithmic}
\end{algorithm}

% =====================================================================
\subsection{Numerical settings}
\label{sec:slpA:settings}

\autoref{tab:slp_settings} lists the numerical values used in the
present study. The worst-case number of MAPDL evaluations per outer
iteration is
\begin{equation}
\label{eq:appA_fea_cost}
    1+n+N_{\mathrm{bt}},
\end{equation}
where the first evaluation is the baseline analysis at $x_k$, the
next $n$ evaluations are the forward finite-difference perturbations,
and the final $N_{\mathrm{bt}}$ evaluations are the maximum number of
inner trial-point evaluations.

\begin{table}[h]
\centering
\small
\begin{tabular}{lll}
\toprule
Parameter & Value & Description \\
\midrule
$\delta_0$                         & $0.10$        & Initial move-limit fraction \\
$\delta_{\max}$                    & $0.10$        & Maximum move-limit fraction \\
$\delta_{\min}$                    & $10^{-5}$     & Minimum move-limit fraction \\
$a_{\mathrm{exp}}$                 & $1.4$         & Move-limit expansion factor \\
$a_{\mathrm{shr}}$                 & $0.5$         & Move-limit shrink factor \\
$R_0$                              & $10^{2}$      & Initial slack penalty \\
$a_M$                              & $2$           & Slack-penalty growth factor \\
$R_{\max}$                         & $10^{6}$      & Maximum slack penalty \\
$\beta$                            & $1-10^{-6}$   & Filter violation parameter \\
$\gamma$                           & $10^{-6}$     & Filter merit parameter \\
$\sigma$                           & $10^{-4}$     & Sufficient-decrease parameter \\
$\gamma_v$                         & $10^{-3}$     & Feasibility-restoration parameter \\
$N_{\mathrm{bt}}$                  & $8$           & Maximum inner LP attempts \\
$\eta_{\mathrm{bt}}$               & $0.5$         & Trust-region backtracking shrink factor \\
$\varepsilon_{\mathrm{feas}}$      & $10^{-5}$     & Feasibility tolerance \\
$\varepsilon_{\mathrm{slack}}$     & $10^{-8}$     & Slack tolerance \\
$\varepsilon_x$                    & $10^{-4}$     & Step-length tolerance \\
$\varepsilon_f$                    & $10^{-3}$     & Merit-change tolerance \\
$\varepsilon_g$                    & $10^{-6}$     & Objective-gradient tolerance \\
$h_{\mathrm{rel}}$                 & $10^{-3}$     & Relative finite-difference step \\
\bottomrule
\end{tabular}
\caption{Numerical parameters used for the LT1292 mast optimization.}
\label{tab:slp_settings}
\end{table}

% =====================================================================
\begin{thebibliography}{9}

\bibitem{sorensen2026fminslp}
S{\o}rensen, R.,
\textit{fminslp: A Sequential Linear Programming Framework with a
Global Convergence Filter}.
MATLAB implementation, Aalborg University, 2026.

\bibitem{sorensen2015thicknessfilters}
S{\o}rensen, R.\ and Lund, E.,
``Thickness filters for gradient-based multi-material and thickness
optimization of laminated composite structures,''
\textit{Structural and Multidisciplinary Optimization},
vol.~52, no.~2, pp.~227--250, 2015.
doi:10.1007/s00158-015-1230-3.

\end{thebibliography}

\end{document}